\section{The complete conductor graph: a quantitative full-volume comparison}\label{the-complete-conductor-graph-a-quantitative-full-volume-comparison}

21 September 2026. This is an independent continuation of KF1--55 and OR1--25 on their original simple-quartet domain. The actual conductor, period, root grid, invariant columns and source order remain fixed.

The new evaluated assertion is CG24 below: the metric shear changes the determinant of the \textbf{entire \(g\)-dimensional conductor quotient} by \(O_A(q\log(q+2))=o(kq)\), at every original cutoff. This follows from a full polynomial-norm estimate that retains the exceptional conductor roots, a quantitative angle estimate and the complete WCF exterior bound. It does not yet bound the shear on the specified proper subspace \(\overline K\). CG25--28 give the exact residual expression on that unchanged subspace; the full-volume estimate cannot be substituted for it.

\subsection{1. Original objects and finite domain}\label{original-objects-and-finite-domain}

Keep all notation and coefficient maps of KF1--55. In particular \[
q=(k+1)^2,\quad q'=(k-7)^2,\quad
v=\operatorname{ord}_0E_A\le80,\quad
g=q-q'-v,\quad a_A=\mu_v\binom qv\ne0,
\tag{CG1}
\] and \[
\mathcal U_Ah=\sum_{r,s=0}^8a_{rs}d_{rs}(S')h(S'+b_{rs}),
\quad d_A=\mathcal U_A1,\quad
\chi'\mathcal U_Ah=\mathcal T_A(\chi h).
\tag{CG2}
\] Here \(\chi'\) names the lower root polynomial, not a derivative. The entire relation and target metric are \[
Y_N=\mathcal P'_{L+g}/\mathcal U_A\mathcal P_L,
\quad L=N-q\in\{-1,0,q-1,q\},
\quad\|f\|_H^2=\int_{\mathbb R}|\chi'(c'+iy)f(c'+iy)|^2d\sigma(y),
\tag{CG3}
\] where \(c'=k/2-4\) and \[
d\sigma(y)=\frac{|\Gamma(1/4+iy/2)|^2}{2\pi}dy,
\quad M_\sigma=\sqrt{2\pi},
\quad \frac{d\sigma}{dy}\ge C_Le^{-\pi|y|},
\quad C_L=\frac\pi{\Gamma(3/4)^2}.
\tag{CG4}
\] The mass, recurrence and Gamma lower bound are the established original ones. The lower bound also follows directly from Gamma reflection and the Euler integral bound on \(\Gamma(3/4-iy/2)\), as in OSP11.

Put \(D_0=\sqrt{\delta^2+\gamma^2}\), \(R_k=kD_0\), \(B_k=4+R_k\), and retain \(r_*,a_*,A_1,B\) exactly from OR4, OR7 and OR14. In particular every zero \(\zeta\) of \(d_A(c'+iy)\) satisfies \(|\zeta|<q/r_*\). Fix the numerical constant \[
\epsilon=2^{-64},\qquad T=\epsilon q.
\tag{CG5}
\] The following estimates use the finite eventual domain \[
k\ge29,\quad q\ge2v,\quad k\ge2v-1,\quad
B_k\le T/8,\quad q\ge8/\epsilon,
\quad u_4:=8/T+(16+R_k^2)/T^2<1/2.
\tag{CG6}
\] For each fixed original conductor these conditions hold eventually. No statement at an excluded small cutoff is inferred by asymptotic notation. The exact identities below also apply at those cutoffs.

\subsection{2. Every Gamma linear factor has a polynomial lower bound}\label{every-gamma-linear-factor-has-a-polynomial-lower-bound}

The original monic polynomials and norms are \[
p_0=1,\ p_1=y,\ p_{j+1}=yp_j-j(j-1/2)p_{j-1},
\quad \gamma_j=M_\sigma j!(1/2)_j,
\quad \phi_j=p_j/\sqrt{\gamma_j}.
\tag{CG7}
\] Their exact generating function is \[
\sum_{j\ge0}p_j(y)t^j/j!=(1+t^2)^{-1/4}e^{y\arctan t}.
\tag{CG8}
\] This is the Meixner--Pollaczek formula with \(p_j(y)=j!P_j^{(1/4)}(y/2;\pi/2)\); the primary source is \href{https://dlmf.nist.gov/18.23.E7}{DLMF 18.23.7}, with recurrence \href{https://dlmf.nist.gov/18.22.E8}{18.22.8}, by T. H. Koornwinder, R. Wong, R. Koekoek, R. F. Swarttouw and W. P. Reinhardt. Differentiating this analytic identity at \(t=0\) gives the finite formula \[
\partial_y\phi_n
=\sum_{\substack{1\le a\le n\\a\ \mathrm{odd}}}
\frac{(-1)^{(a-1)/2}}a\frac{\rho_n}{\rho_{n-a}}\phi_{n-a},
\quad\rho_n=\sqrt{n!/(1/2)_n}.
\tag{CG9}
\] The condition in this sum is simply that \(a\) is odd. Since \(1\le\rho_n\le\sqrt{2n+1}\), every row sum and every column sum of its matrix on degree at most \(D\) is bounded by \[
b_D=\sqrt{2D+1}\,H_D,\qquad H_D=\sum_{j=1}^D1/j,
\quad b_0=0.
\tag{CG10}
\] For any matrix \(A\), Cauchy--Schwarz in each row with weights \(|A_{ij}|\), followed by summing columns, proves \(\|A\|_2^2\le\|A\|_1\|A\|_\infty\). Thus \(\|\partial_y\|_{\mathcal P_D\to\mathcal P_{D-1}}\le b_D\).

For every complex \(z\) and polynomial \(p\) of degree at most \(D\), \[
\boxed{\frac{\|p\|_\sigma}{2b_{D+1}}
\le\|(y-z)p\|_\sigma
\le[2(D+1)+|z|]\|p\|_\sigma.}
\tag{CG11}
\] For the lower bound use \(p=\partial_y((y-z)p)-(y-z)\partial_yp\), pair with \(p\), and move the multiplication operator in the second term to the other side. Its adjoint is multiplication by \(y-\bar z\). On the real integration line \(\|(y-z)p\|_\sigma=\|(y-\bar z)p\|_\sigma\). The derivative bounds at degrees \(D+1,D\) give \(\|p\|^2\le(b_{D+1}+b_D)\|p\|\|(y-z)p\|\), proving the first inequality. The upper bound follows from CG7: the raising and lowering coefficients have norms at most \(D+1,D\), respectively. This argument holds for every complex \(z\), including a real zero inside the region carrying the norm. No factor is removed.

For the degrees below define the convenient common bounds \[
\lambda_q=(2b_{2q+1})^{-1},\quad
U_q=4q+2+q/r_*,\quad U_q^0=4q+2.
\tag{CG12}
\] Every multiplication needed below has input degree at most \(2q\). A product of \(b\) factors with zeros of modulus at most \(q/r_*\) therefore has norm between \(\lambda_q^b\) and \(U_q^b\) times the input norm, provided the entire product stays within the displayed degree range. The same bounds apply to \(y^b\). Their comparison costs at most \(2b\log(U_q/\lambda_q)\) in the logarithm of squared norms.

\subsection{3. All conductor roots, including those on the q-scale}\label{all-conductor-roots-including-those-on-the-q-scale}

Split the actual root list, with multiplicities, at \(T/2\): \[
d_A(c'+iy)=a_Ai^g D_{\rm in}(y)D_{\rm out}(y),
\quad D_{\rm in}=\prod_{|\zeta|\le T/2}(y-\zeta),
\quad D_{\rm out}=\prod_{|\zeta|>T/2}(y-\zeta).
\tag{CG13}
\] Write \(b=\deg D_{\rm out}\), \(a=g-b\). OR10--11 applied with reciprocal radius \(2/\epsilon\) is allowed by CG6 and proves \[
b\le B_A:=\frac{\log^+(A_1/a_*)+8B/\epsilon}{\log2}.
\tag{CG14}
\] Thus \(b\) has a fixed finite bound for the original conductor. This is stronger than counting roots outside the shrinking relative disk in OR13. It has a different, explicitly fixed, radius. Retain also OR16's full root sum \(S_k=\sum_\zeta|\zeta|=O_A(q\log(q+2))\).

For every polynomial \(p\) of degree \(d\le2q\) and every \(Q\ge q/2\), with \(Q+a+d\le2q\), one has \[
e^{-E_{\rm in}(Q,d)}\|y^{Q+a}p\|_\sigma^2
\le\|y^QD_{\rm in}p\|_\sigma^2
\le e^{E_{\rm in}(Q,d)}\|y^{Q+a}p\|_\sigma^2,
\tag{CG15}
\] where the completely finite error is \[
\begin{split}
\eta_k&=4S_k/T,\\
K(Q,a,d)&=\frac{M_\sigma}{C_L}\frac{(d+1)^2}{q}
 e^{2\pi q}10^{2d}\epsilon^{2Q}(3\epsilon/2)^{2a},\\
E_{\rm in}(Q,d)&=\eta_k+\log(1+K(Q,a,d)).
\end{split}\tag{CG16}
\] Here \(E_{\rm in}=O_A(\log(q+2))\) uniformly in these degrees.

For completeness, outside \((-T,T)\), each root satisfies \(|\zeta/y|\le1/2\), so \[
\left|\log\frac{|D_{\rm in}(y)|^2}{|y|^{2a}}\right|
\le 2\sum\frac{|\zeta/y|}{1-|\zeta/y|}\le4S_k/T.
\] To retain the inner interval use the following finite polynomial bound: \[
\sup_{|y|\le T}|p(y)|^2
\le\frac{(d+1)^2 10^{2d}}q\int_q^{2q}|p(u)|^2du.
\tag{CG17}
\] Indeed expand \(p(qz)\) in the orthonormal polynomials \(\sqrt{2j+1}L_j(2z-3)\) on \([1,2]\). Rodrigues' formula and \(j\) integrations by parts give \(\int_{-1}^1L_j^2=2/(2j+1)\). The finite formula \(L_j(x)=2^{-j}\sum_{h=0}^j\binom jh^2(x-1)^{j-h}(x+1)^h\) gives \(|L_j(x)|\le10^j\) for \(|x|\le4\), using \(\sum_h\binom jh^2\le4^j\). Since \(|2z-3|\le4\) on the required inner interval, Cauchy--Schwarz proves CG17. On that interval \(|y|^{2Q}|D_{\rm in}(y)|^2\le T^{2Q}(3T/2)^{2a}\). The integral over \([q,2q]\) of the reference norm is at least \(q^{2(Q+a)}C_Le^{-2\pi q}\int_q^{2q}|p|^2\). CG17 and the full mass \(M_\sigma\) therefore bound both inner contributions by \(K(Q,a,d)\) times that outside contribution. Combining this with the outside logarithmic comparison gives CG15. Finally \(Q\ge q/2,d\le2q,3\epsilon/2<1\) bound \(K\) by \(M_\sigma(2q+1)^2(C_Lq)^{-1}
 \exp\{q(2\pi+4\log10+\log\epsilon)\}\), which decays exponentially. This proves the stated order without losing the inner region or a root lying on the real axis.

The original lower polynomial satisfies the same complete estimate \[
e^{-E_{\chi'}}\|y^{q'}p\|_\sigma^2
\le\|\chi'(c'+iy)p(y)\|_\sigma^2
\le e^{E_{\chi'}}\|y^{q'}p\|_\sigma^2,
\quad E_{\chi'}=O_A(1),
\tag{CG18}
\] uniformly for degree \(p\le2q-q'\). Here and below modulus-one leading phases do not change the norms. An explicit bound, also needed for the shifted upper polynomial, is as follows. For \(w\in\{0,4\}\), \(Q\in\{q',q\}\), put \[
\begin{split}
u_w&=2w/T+(w^2+R_k^2)/T^2,\qquad
\eta_{Q,w}=Qu_w/(1-u_w),\\
K_{Q,w,d}&=\frac{M_\sigma}{C_L}\frac{(d+1)^2}{q}
e^{2\pi q}10^{2d}
\left(\frac{(T+w)^2+R_k^2}{q^2}\right)^Q,\\
E_{Q,w,d}&=\eta_{Q,w}+\log(1+K_{Q,w,d}).
\end{split}\tag{CG19}
\] Use \(E_{\chi'}=E_{q',0,2q-q'}\). The roots of either centered original polynomial pair as \(\omega,-\omega\). For the polynomial shifted by \(iw\), each paired factor divided by \(y^2\) is \(1+2iw/y-(w^2+\omega^2)/y^2\). Its deviation from one is at most \(u_w<1/2\). This proves the outside logarithmic bound \(\eta_{Q,w}\). Inside the interval its squared modulus is at most \(((T+w)^2+R_k^2)^Q\), so the same proof of CG17--16 supplies \(K_{Q,w,d}\). The guards imply that \(((T+w)^2+R_k^2)/q^2\le4\epsilon^2\), making the latter exponentially small for \(Q\ge q/2,d\le2q\). Also \(\eta_{Q,w}=O_A(1)\). This proves CG18 and the corresponding claim for the upper polynomial in the lower physical coordinate.

Apply CG18, then CG15 with \(Q=q'\) and \(p=D_{\rm out}h\), and finally CG11 to every factor of \(D_{\rm out}\) and to every factor of \(y^b\). For \(0\le L\le q\), \(\deg h\le L\), all degrees remain at most \(q'+g+L=q-v+L\le2q\). Thus \[
\left|\log\frac{\|\chi' d_Ah\|_{\sigma,c'}^2}
 {|a_A|^2\|y^{q-v}h(c'+iy)\|_\sigma^2}\right|
\le E_{\chi'}+E_{\rm in}(q',L+b)
+2b\log(U_q/\lambda_q)=:E_d.
\tag{CG20}
\] The ratio is evaluated only for nonzero \(h\), so its denominators are positive. The bound is uniform over every coefficient vector, including its worst direction, and \(E_d=O_A(\log(q+2))\).

\subsection{4. Exact physical centers and complete relation norms}\label{exact-physical-centers-and-complete-relation-norms}

For polynomials of degree at most \(M\), translation by \(4i\) in \(y\) and its inverse have Gamma operator norm at most \[
T_M=e(M+1)(2M+1)^{5/2}\qquad(M\ge1).
\tag{CG21}
\] To prove it, CG8 multiplies under this translation by \(e^{\pm4i\arctan t}=((1+it)/(1-it))^{\pm2}\). On \(|t|=M/(M+1)\) its modulus is at most \((2M+1)^2\). Cauchy's coefficient bound through degree \(M\) costs at most \(e\). In the orthonormal Gamma frame the coefficient on the \(a\)th diagonal is multiplied by \(\rho_n/\rho_{n-a}\le\sqrt{2M+1}\). Sum the at most \(M+1\) diagonal norms. The inverse has the same bound. This explicitly carries the original upper center \(c=c'+4\) to the lower center rather than identifying their coefficient norms.

Combining CG20, the \(w=4,Q=q\) case of CG19, CG21 for the entire polynomial \(\chi h\), and the \(v\) factors of \(y\) in CG11 gives \[
\boxed{e^{-E_M}\|\chi h\|_{\sigma,c}^2
\le\|\chi'd_Ah\|_{\sigma,c'}^2
\le e^{E_M}\|\chi h\|_{\sigma,c}^2,\qquad E_M=O_A(\log(q+2)).}
\tag{CG22}
\] One explicit common error for all \(0\le L\le q\) is \[
E_M=\max_{0\le L\le q}E_d
+E_{q,4,q}+2\log T_{2q}+2|\log|a_A||
+2v\max\{\log U_q^0,\log(1/\lambda_q)\}.
\] The maximum can be replaced by the same displayed formulas at \(d=q+\lceil B_A\rceil\), or bounded directly by \(d\le2q\) on their actual range. Only the actual \(b\le g\) is used in the formula when the fixed coarse bound exceeds \(g\).

The equality of coefficient vectors \(h(S)\) and \(h(S')\) is retained in CG22; CG21 is what transports their different evaluation lines. The modulus \(|a_A|\), the \(v\) degrees and all exceptional roots were paid for explicitly. No conjectured absence of such roots or assumed nonvanishing on the integration line enters the estimate.

\subsection{5. The full g-dimensional graph volume is subleading}\label{the-full-g-dimensional-graph-volume-is-subleading}

Let \(t=L+1\) for \(L\ge0\), \(N=q+L\), \(D=N-v\), and let \(E_L=\chi\mathcal P_L\subset\mathcal P_N\), with its original center-\(c\) Gamma metric. The degree-\(<v\) space is denoted \(\mathcal L\). It has zero intersection with \(E_L\). Here is a quantitative angle bound that preserves the entire space. For \(v\ge1\), take the actual upper roots whose \(y\)-coordinates are \[
z_j=-\gamma+i(2j+1)\delta,\qquad0\le j<v.
\] They are present by \(k\ge2v-1\), are distinct, and are independent of growing \(k\). Let \(\ell_j(y)\) be their degree-\(<v\) Lagrange polynomials, and set \[
\begin{gathered} C_I=\left(\sum_j\|\ell_j\|_\sigma^2\right)^{1/2},\qquad R_I=\max_j|z_j|,\\
C_E(N)=\left[\frac{ve^2}{M_\sigma}(N+1)^{3/2}
 (2N+1)^{R_I+1}\right]^{1/2},\\
\theta_N=\min\{1,(C_IC_E(N))^{-1}\}.\end{gathered}
\tag{CG23}
\] Set \(\theta_N=1\) if \(v=0\). The coefficient estimate from CG8 on radius \(N/(N+1)\), with \(|\arctan t|\le\tfrac12\log(2N+1)\), proves \(|\phi_j(z)|^2\le e^2M_\sigma^{-1}(N+1)^{1/2}
 (2N+1)^{|z|+1}\) for \(j\le N\), exactly as in KF33. Summing gives the evaluation norm \(C_E(N)\). Since every element of \(E_L\) vanishes at all \(z_j\), interpolation gives, for \(l\in\mathcal L\), \(\|l\|\le C_I\|\operatorname{ev}l\|
\le C_IC_E(N)\operatorname{dist}(l,E_L)\). Thus this angle has a polynomial lower bound. By taking the adjoint of the orthogonal cross projection, the projection of \(E_L\) onto \(\mathcal L^\perp\) has all singular values at least \(\theta_N\) and at most \(v\) different from one. Its squared determinant is therefore between \(\theta_N^{2v}\) and one.

Apply the complete WCF4--13 map to that quotient. Write \(M_D=\max\{1,B_D^+\}\) and \(J_D=(D+1)\log^+(1/|(-i)^v\mu_v/v!|)\), exactly as in KF46. The forward exterior norm at rank \(t\) is at most \(M_D^t\); the inverse exterior norm is at most \(M_D^{D+1-t}e^{J_D}\). Consequently, in the \textbf{same source coefficient frame \(h\)}, \[
-2(D+1-t)\log M_D-2J_D+2v\log\theta_N
\le\log\frac{\det\operatorname{Gram}(\chi'\mathcal U_Ah)}
{\det\operatorname{Gram}(\chi h)}
\le2t\log M_D.
\] Together with CG22, this gives the finite common bound \[
\left|\log\frac{\det\operatorname{Gram}(\chi'\mathcal U_Ah)}
{\det\operatorname{Gram}(\chi'd_Ah)}\right|
\le B_N:=2(N+1)\log M_D+2J_D+2v\log(1/\theta_N)+tE_M.
\] It is \(O_A(q\log(q+2))\). The inverse bound is an exterior bound before multiplying by a rank, so it does not incur a spurious \(q^2\) cost. The polynomial angle was proved above before its fixed rank was used.

Let \(Q_N^U\) be the complete \(g\)-dimensional quotient Gram CG3 in the unchanged low remainder frame \(1,S',\ldots,S'^{g-1}\), and let \(Q_N^d\) be the quotient by \(d_A\mathcal P_L\) in that same frame with the \textbf{untransformed} source metric \(H\). KF52 gives the full coefficient matrix \(\mathcal C_L\). The change of basis from \([1,\ldots,S'^{g-1},d_A,d_AS',\ldots,d_AS'^L]\) to the basis with \(\mathcal U_A1,\ldots,\mathcal U_AS'^L\) in its last slots has determinant \(\det\mathcal C_L\). Block Gram determinants, with every cross block retained, therefore give \[
\begin{split}
\log\det Q_N^U-\log\det Q_N^d
={}&2\log\det\mathcal C_L
-\log\frac{\det\operatorname{Gram}(\chi'\mathcal U_Ah)}
{\det\operatorname{Gram}(\chi'd_Ah)},\\
\boxed{\begin{gathered}|\log\det Q_N^U-\log\det Q_N^d|
\le 2v(L+1)\log\frac{q+L}{q-v+1}+B_N,\\ B_N=O_A(q\log(q+2))=o(kq).\end{gathered}}
\end{split}\tag{CG24}
\] At \(L=-1\) both quotients are the identical whole low space. At \(L=0\), \(\mathcal U_A1=d_A\), \(\mathcal C_0=1\), \(\mathcal R_0=0\), and the quotients are again identical. Thus the metric shear is exactly the identity at \textbf{both low cutoffs}, and only the two high errors occur in its four-return. CG24 evaluates their total full-target determinant cost; it does not use an assumed isometry of the untransformed metrics.

\subsection{6. The exact unresolved receiver on the original kernel plane}\label{the-exact-unresolved-receiver-on-the-original-kernel-plane}

Write the actual complete source Gram in the polynomial-division basis \(r+d_Aw\) as \[
H=\begin{pmatrix}A&B\\B^*&D\end{pmatrix},\quad
F=\mathcal R_L\mathcal C_L^{-1},\quad
Q_N^d=A-BD^{-1}B^*.
\tag{CG25}
\] Its relation graph is \((Fw,w)\); its attained Gram is therefore \[
Q_N^U=A-(AF+B)
(D+B^*F+F^*B+F^*AF)^{-1}(F^*A+B^*).
\tag{CG26}
\] This is obtained by expanding and minimizing \(\langle(r+Fw,w),H(r+Fw,w)\rangle\) in \(w\). Every inverse is of a positive definite Gram of an injective relation map. Equivalently let \(R_d,D_d\) be the remainder and quotient-coefficient maps of ordinary polynomial division by \(d_A\). The actual quotient map is exactly \(R_d-FD_d\), since it kills \((Fw,w)\) and fixes every low representative. Consequently \[
(Q_N^U)^{-1}=(R_d-FD_d)H^{-1}(R_d-FD_d)^*,
\quad (Q_N^d)^{-1}=R_dH^{-1}R_d^*.
\tag{CG27}
\] These formulas retain the transformed metric \((\mathcal J_L^{-1})^*H\mathcal J_L^{-1}\) from KF55: CG26--27 are its quotient minima in the original low frame, not a replacement of that transformed metric.

Let \(E_K\) be the exact low coefficient frame of \(\overline K=\Psi_N(K/(K\cap\mathcal L))\), defined by KF7 and KF10 with the complete \(Z=[M_A,S_kZ_k]\). It is independent of \(N\). The still unevaluated signed graph effect in the requested allocation is \[
\mathcal G_K=
-\sum_{N\in\{2q-1,2q\}}
\log\frac{\det(E_K^*Q_N^UE_K)}{\det(E_K^*Q_N^dE_K)}.
\tag{CG28}
\] The missing low terms are exactly zero by CG24, not discarded by an estimate. Formula CG24 controls the corresponding full \(g\)-volume. It does not control CG28 by its rank. For example positive matrices \(\operatorname{diag}(e^a,e^{-a})\) and \(I_2\) have identical determinants but their restrictions to the first coordinate differ by \(a\). This elementary example is only a counterexample to that inference, not a replacement for the original matrices. The original connecting maps and exact remaining scalar are CG25--28.

\subsection{7. The remaining allocation is an explicit original invariant receiver}\label{the-remaining-allocation-is-an-explicit-original-invariant-receiver}

The residual in CG28 can now be transferred, with the proved subleading full-volume cost, to the complete remaining invariant observations. This identifies the next matrix without changing their rank or columns. Use exactly KF8--10, and write their coefficient matrices as \[
B_{\mathcal L}=J_R^{\mathsf T}V\big|_{\mathcal L},\qquad
B_\eta\eta=J_R^{\mathsf T}V\mathcal R(\chi'\eta),
\quad J_R=S_kZ_k.
\tag{CG29}
\] Here \(\mathcal R\) is KF8's fixed right inverse, not a cutoff return. Let \(\Pi\) be the Euclidean orthogonal projection onto the complement of \(\operatorname{im}B_{\mathcal L}\) in the displayed observation coefficient space. This auxiliary Euclidean projection performs only linear constraint elimination; it does not replace any original Gamma metric. Choose once and for all a row basis of \(\Pi B_\eta\), denoted \(A_R:\mathbb C^g\to\mathbb C^r\). Then \[
\ker A_R=\overline K,\qquad
r=g-(m-s_k)=8k-32-v+s_k.
\tag{CG30}
\] Indeed \(\eta\) lies in the kernel precisely when \(B_\eta\eta\in\operatorname{im}B_{\mathcal L}\), which is equivalent to existence of \(l\in\mathcal L\) solving the complete KF10 constraint. KF7 identifies exactly the quotient by \(K\cap\mathcal L\), proving the kernel and rank assertions. The full conductor columns were used before this map, and every column of \(J_R\) remains in CG29.

For \(X\in\{U,d\}\) put \[
C_{R,N}^X=A_R(Q_N^X)^{-1}A_R^*,\qquad
\Delta_{R,N}=\log\det C_{R,N}^U-\log\det C_{R,N}^d.
\] These are the actual attained quotient covariances of that fixed observation map. If \(r=0\), their empty determinant is one. The exact block determinant identity is \[
\boxed{\log\frac{\det(E_K^*Q_N^UE_K)}{\det(E_K^*Q_N^dE_K)}
=\log\frac{\det Q_N^U}{\det Q_N^d}+\Delta_{R,N}.}
\tag{CG31}
\] To prove it choose a fixed right inverse \(T_R\) of \(A_R\) and form the invertible frame \(P=[E_K,T_R]\). The Schur complement of the first block of \(P^*Q_N^XP\) is the attained quotient metric \((A_R(Q_N^X)^{-1}A_R^*)^{-1}\), since \(A_RE_K=0\) and \(A_RT_R=I_r\). Taking determinants gives \(|\det P|^2\det Q_N^X
=\det(E_K^*Q_N^XE_K)/\det C_{R,N}^X\). The fixed frame determinant cancels between \(U,d\), proving CG31.

In particular CG24 gives the proved finite receiver \[
\left|\mathcal G_K+
\Delta_{R,2q-1}+\Delta_{R,2q}\right|
\le\sum_{L\in\{q-1,q\}}
\left[2v(L+1)\log\frac{q+L}{q-v+1}+B_{q+L}\right]
=o_A(kq).
\tag{CG32}
\] Thus the full shear-volume calculation removes that entire part of the allocation at its required scale. The remaining leading part is the explicit signed sum of the two original invariant covariances in CG32. Its matrix is obtained by inserting \(A_R\) on both sides of CG27; its actual coefficients, all cross terms and attained minima remain. The present proof does not yet evaluate those two logarithms.

\subsection{Proof-source coverage and citations}\label{proof-source-coverage-and-citations}

\begin{itemize}
\tightlist
\item
  KF1--55, \emph{The actual invariant kernel: degree filtration, quantitative minors, and its complete conductor graph}, read completely. Its complete source must accompany publication of this continuation until a pinned public proof-file URL has been verified.
\item
  OR1--25, \emph{The roots of the complete conductor polynomial}, read completely. CG14 and CG20 use OR10--16, including every finite guard. Its complete source must likewise accompany publication until its public proof-file URL is verified; no live URL is invented here.
\item
  WCF4--13 and WCF19a--f, \emph{A polynomial forward bound for the full weighted conductor}, \href{https://github.com/KokunoYumeto/zeta-function-research-reader/blob/52274804051e325c8a6ef69c17b3958cc2caba0a/workbenches/splitzero-tandem/continuations/20260920-gamma-and-finite-derivative/providers/WEIGHTED_CONDUCTOR_FORWARD.tex}{pinned complete proof file}. The local source was read through WCF20; the exact WCF12--13 and WCF19a--f exterior statements were read directly and used as the accepted provider. No independent full rereading of the remainder of WCF is claimed. The published path and commit are identified by its existing public-readback receipt.
\item
  OSP10--26, in the prior complete source edition, read in full. CG17 and CG19 reproduce the finite estimates needed here, retaining the original Legendre proof, Gamma mass and both integration regions. No unsupported public OSP link is asserted.
\item
  NIST DLMF chapter 18, version 1.2.8, released 15 September 2026, human authors named at CG8. The official 18.23.7 formula and 18.22.8 recurrence were checked, and their original NIST TeX encodings were fetched, read and preserved in \nolinkurl{graph_primary_sources/}. This argument derives its derivative estimate from those formulas; it does not attribute CG11 or CG24 to DLMF.
\end{itemize}

No subleading bound on CG28, no individual projected-current sign and no RH conclusion is claimed by the present derivation.

The pinned WCF public source was downloaded and compared byte for byte with the local source on 21 September 2026; they agree, with SHA256 \nolinkurl{d86dcb3daf13c9b2d5a0aaa4d72cce35a1aff48066ad69c7d9846b9d097dc354}. The independent exact checker \nolinkurl{check_graph_metric.py} passed 81 tests of CG9--11, CG24--27 and CG31. It includes a negative control showing why an equal full determinant cannot assign the determinant on the first fixed direction. The checks concern exact identities and finite auxiliary Grams; the asymptotic estimates are proved in the text above.
